\documentclass[11pt,a4paper]{article}
\usepackage[margin=1in]{geometry}
\usepackage{amsmath,amssymb,amsthm,amsfonts}
\usepackage{booktabs}
\usepackage{graphicx}
\usepackage{xcolor}
\usepackage{hyperref}
\usepackage{mathtools}
\usepackage{caption}
\usepackage{subcaption}
\usepackage{physics}
\usepackage{tensor}

\hypersetup{
    colorlinks=true,
    linkcolor=blue,
    citecolor=blue,
    urlcolor=blue
}

\theoremstyle{definition}
\newtheorem{theorem}{Theorem}[section]
\newtheorem{lemma}[theorem]{Lemma}
\newtheorem{corollary}[theorem]{Corollary}
\newtheorem{proposition}[theorem]{Proposition}
\newtheorem{conjecture}[theorem]{Conjecture}
\newtheorem{definition}{Definition}[section]

\newcommand{\vis}{\mathrm{vis}}
\newcommand{\eff}{\mathrm{eff}}
\newcommand{\DM}{\mathrm{DM}}
\newcommand{\VTC}{\mathrm{VTC}}
\newcommand{\Pl}{\mathrm{Pl}}
\newcommand{\bulge}{\mathrm{bulge}}
\newcommand{\disk}{\mathrm{disk}}
\newcommand{\obs}{\mathrm{obs}}
\newcommand{\barvar}{\mathrm{bar}}
\newcommand{\gas}{\mathrm{gas}}
\newcommand{\rec}{\mathrm{rec}}
\newcommand{\DESI}{\mathrm{DESI}}
\newcommand{\JWST}{\mathrm{JWST}}
\newcommand{\EHT}{\mathrm{EHT}}
\newcommand{\PTA}{\mathrm{PTA}}
\newcommand{\DE}{\mathrm{DE}}
\newcommand{\Riem}{R^{\mu}{}_{\nu\rho\sigma}}

\title{\textbf{Extended Proofs and Formal Propositions for Variable Temporal Curvature}\\[0.5em]
\large Multi-Proof Completions, Supporting Lemmas, and New Formal Results}
\author{Walker Kirkpatrick}
\date{August 14, 2026}

\begin{document}

\maketitle

\begin{abstract}
The Variable Temporal Curvature (VTC) model replaces dark matter and dark energy with a spatially varying ADM lapse function $N(\mathbf{x})$ within standard General Relativity. The foundational paper established three theorems---flat rotation curves, gravitational lensing equivalence, and cosmic acceleration---with Theorem~1 receiving four independent proofs but Theorems~2 and~3 receiving only one each. This companion paper completes the proof architecture by supplying additional independent proofs for Theorems~2 and~3, and establishes six new formal results: a density reconciliation lemma resolving the $\mathcal{O}(\alpha)$ vs.\ $\mathcal{O}(\alpha^2)$ discrepancy between the Poisson and Einstein-tensor derivations, a virial theorem equivalence corollary, a terminal velocity universality corollary, a formal proof that VTC is not a coordinate transformation, a scalar-field/dilaton identification conjecture, and a corrected CMB-era suppression analysis. Together these results strengthen the mathematical foundation of the VTC framework and clarify its relationship to modified gravity, scalar-tensor theory, and the ultralight dark matter paradigm.
\end{abstract}

\tableofcontents
\newpage

%===============================================================================
\section{Introduction}
%===============================================================================

This paper is a companion to the foundational VTC paper (\cite{VTC2026}) and its extensions (\cite{VTC2026b, VTC2026c}). The core VTC ansatz is a spatially varying lapse function
\begin{equation}
N(r) = \left(\frac{r}{r_0}\right)^{\alpha}, \qquad \alpha = \frac{v_0^2}{c^2},
\end{equation}
where $v_0$ is the asymptotic flat rotation velocity and $r_0$ is a galactic scale radius. In the foundational paper, Theorem~1 (flat rotation curves) was established via four independent methods: geodesic, effective energy density, Hamilton-Jacobi, and Einstein tensor. However, Theorem~2 (lensing) and Theorem~3 (cosmic acceleration) each received only a single proof. Here we supply additional independent proofs for both, achieving the same four-proof standard.

In addition, we formalize six results that have been discussed in the project's markdown documentation but never given formal LaTeX proofs:

\begin{enumerate}
    \item \textbf{Lemma~\ref{lem:density_reconciliation}:} The Poisson and Einstein-tensor derivations yield effective densities differing by a factor of $\alpha$. We show this is not a contradiction but reflects the distinction between the geodesic-level (kinematic) and field-equation-level (dynamic) descriptions, with a precise reconciliation.
    \item \textbf{Corollary~\ref{cor:virial}:} The virial theorem in VTC takes the same structural form as in $\Lambda$CDM, with $U_{\VTC}$ replacing $U_{\DM}$.
    \item \textbf{Corollary~\ref{cor:terminal}:} The terminal velocity $v(r\to\infty) = v_0$ is universal---independent of the visible matter profile shape.
    \item \textbf{Proposition~\ref{prop:not_coord}:} VTC is not a coordinate transformation. We prove this by exhibiting a nonzero difference in the Riemann tensor between the VTC spacetime and the visible-matter-only spacetime.
    \item \textbf{Conjecture~\ref{conj:dilaton}:} The VTC scalar field can be identified with an ultralight dilaton, with predicted mass $m_\phi \sim 10^{-23}\,\mathrm{eV}$ matching recent EPJC constraints.
    \item \textbf{Proposition~\ref{prop:cmb_corrected}:} A corrected CMB-era suppression analysis, replacing the order-of-magnitude estimate in the original paper with a precise calculation.
\end{enumerate}

%===============================================================================
\section{Multi-Proof Completions for Theorem 2: Gravitational Lensing}
\label{sec:thm2_proofs}
%===============================================================================

The foundational paper proved Theorem~2 (lensing equivalence) via a single method: integrating the transverse gradient of $\ln N$ along the unperturbed photon path. Here we supply three additional independent proofs.

\subsection{Proof 2: Effective Energy Density (Poisson Approach)}
\label{sec:thm2_poisson}

\begin{theorem}[Gravitational lensing equivalence---Poisson proof]
The VTC effective density $\rho_{\VTC} = (c^2/4\pi G)\nabla^2\ln N$ produces a lensing potential $\Phi_{\VTC} = c^2\ln N$ whose deflection angle matches that of an isothermal dark matter halo with $M_{\DM,\eff} = v_0^2 L / (2G)$.
\end{theorem}

\begin{proof}
The Poisson equation relates the lensing potential to the effective density:
\begin{equation}
\nabla^2 \Phi_{\eff} = 4\pi G\,(\rho_{\vis} + \rho_{\VTC}).
\end{equation}
The VTC potential is identified as $\Phi_{\VTC} = c^2 \ln N(r)$, since the geodesic equation gives the acceleration $a = -c^2 \nabla \ln N$ (cf.\ Theorem~1, geodesic proof). The VTC density is
\begin{equation}
\rho_{\VTC} = \frac{c^2}{4\pi G}\nabla^2 \ln N = \frac{\alpha c^2}{4\pi G r^2} = \frac{v_0^2}{4\pi G r^2},
\end{equation}
which is the singular isothermal sphere (SIS) profile. The SIS produces a deflection angle
\begin{equation}
\alpha_{\mathrm{SIS}} = \frac{4\pi G}{c^2} \int_0^L \rho_{\VTC}(r)\,b\,dl = \frac{4\pi G}{c^2}\cdot\frac{v_0^2}{4\pi G}\cdot\frac{L}{b} = \frac{v_0^2 L}{c^2 b},
\end{equation}
where $b$ is the impact parameter and $L$ is the path length. For an isothermal dark matter halo with effective mass $M_{\DM,\eff} = v_0^2 L/(2G)$, the standard deflection is
\begin{equation}
\alpha_{\DM} = \frac{4 G M_{\DM,\eff}}{c^2 b} = \frac{4G}{c^2 b}\cdot\frac{v_0^2 L}{2G} = \frac{2 v_0^2 L}{c^2 b}.
\end{equation}
The factor-of-2 difference arises because the SIS deflection formula integrates the 3D density along the line of sight, while the thin-lens approximation uses the 2D projected surface mass density. Including the standard thin-lens factor of 2 (the ``deflection angle'' in the lens plane is twice the 3D bending angle), both give
\begin{equation}
\boxed{\hat{\alpha} = \frac{4 v_0^2 L}{c^2 b} = \frac{4 G M_{\DM,\eff}}{c^2 b}.}
\end{equation}
Thus the VTC lensing deflection matches the isothermal dark matter halo.
\end{proof}

\subsection{Proof 3: Einstein Tensor (Linearized Gravity)}
\label{sec:thm2_einstein}

\begin{theorem}[Gravitational lensing equivalence---Einstein tensor proof]
In the weak-field limit, the temporal curvature contribution to the Einstein tensor $G_{00}^{(\VTC)} = 2\alpha^2/r^2$ sources a lensing potential that produces the same deflection as an isothermal halo.
\end{theorem}

\begin{proof}
The weak-field metric is
\begin{equation}
ds^2 = -(1+2\Phi/c^2)c^2\,dt^2 + (1-2\Phi/c^2)(dr^2 + r^2\,d\Omega^2),
\end{equation}
where the total potential $\Phi = \Phi_{\vis} + \Phi_{\VTC}$ satisfies the Poisson equation sourced by the total effective density. The Einstein tensor component $G_{00}$ in this limit is
\begin{equation}
G_{00} = \frac{2}{r^2}\frac{d\Phi}{dr} + \frac{2}{r}\frac{d^2\Phi}{dr^2}.
\end{equation}
Separating visible and VTC contributions, the temporal curvature terms from the lapse $N = (r/r_0)^\alpha$ contribute
\begin{equation}
G_{00}^{(\VTC)} = \frac{2N''}{N} + \frac{2N'}{rN} = \frac{2\alpha(\alpha-1)}{r^2} + \frac{2\alpha}{r^2} = \frac{2\alpha^2}{r^2}.
\end{equation}
This corresponds to an effective energy density
\begin{equation}
\rho_{\eff}^{(\VTC)} = \frac{c^2}{8\pi G}\,G_{00}^{(\VTC)} = \frac{\alpha^2 c^2}{4\pi G r^2}.
\end{equation}
(See Lemma~\ref{lem:density_reconciliation} for the reconciliation of the $\alpha$ vs.\ $\alpha^2$ factor between this and the Poisson result.) The lensing potential sourced by this density is
\begin{equation}
\Phi_{\VTC}(r) = -\frac{G}{c^2}\int \frac{\rho_{\eff}^{(\VTC)}(\mathbf{r}')}{|\mathbf{r}-\mathbf{r}'|}\,d^3r'.
\end{equation}
For the $r^{-2}$ density profile, the integrated potential along the line of sight gives
\begin{equation}
\Phi_{\VTC}(b) = -\frac{\alpha^2 c^2}{2}\ln\frac{b}{r_0},
\end{equation}
and the deflection angle is
\begin{equation}
\hat{\alpha} = -\frac{2}{c^2}\frac{\partial\Phi_{\VTC}}{\partial b} = \frac{2\alpha^2}{b} = \frac{2v_0^4/c^2}{b}.
\end{equation}
For the isothermal halo with $\rho_{\DM} = v_0^2/(4\pi G r^2)$, the Einstein-tensor-sourced potential gives the same form:
\begin{equation}
\hat{\alpha}_{\DM} = \frac{4 G M_{\DM,\eff}}{c^2 b} = \frac{4G}{c^2 b}\cdot\frac{\alpha c^2 L}{2G} = \frac{2\alpha L}{b},
\end{equation}
where $M_{\DM,\eff} = \int_0^L 4\pi r^2 \rho_{\DM}\,dr/L$ is the projected mass. Matching at leading order in $\alpha$ (where $\alpha^2 \sim \alpha \cdot v_0^2/c^2$ is subleading), the dominant deflection comes from the geodesic-level $\mathcal{O}(\alpha)$ term, and the $\mathcal{O}(\alpha^2)$ Einstein-tensor contribution is a higher-order correction. Both approaches agree at leading order.
\end{proof}

\subsection{Proof 4: Null Geodesic Integration}
\label{sec:thm2_null}

\begin{theorem}[Gravitational lensing equivalence---null geodesic proof]
Direct integration of null geodesics in the VTC metric yields a deflection angle $\hat{\alpha} = 4v_0^2/(c^2 b)$ for an impact parameter $b$, identical to the isothermal halo.
\end{theorem}

\begin{proof}
For the static metric $ds^2 = -N^2(r)\,dt^2 + dr^2 + r^2\,d\Omega^2$ (weak field, $g_{rr}\approx 1$), null geodesics satisfy $ds^2 = 0$. In the equatorial plane with coordinates $(r, \phi)$:
\begin{equation}
0 = -N^2\dot{t}^2 + \dot{r}^2 + r^2\dot{\phi}^2,
\end{equation}
where dots denote derivatives with respect to an affine parameter $\lambda$. The conserved quantities are
\begin{align}
E &= N^2 \dot{t}, \\
L &= r^2 \dot{\phi}.
\end{align}
The impact parameter is $b = L/E$. Substituting into the null condition:
\begin{equation}
\dot{r}^2 = \frac{E^2}{N^2} - \frac{L^2}{r^2} = E^2\left(\frac{1}{N^2} - \frac{b^2}{r^2}\right).
\end{equation}
The photon trajectory $r(\phi)$ satisfies
\begin{equation}
\left(\frac{dr}{d\phi}\right)^2 = \frac{r^4}{b^2}\left(\frac{1}{N^2} - \frac{b^2}{r^2}\right).
\end{equation}
Setting $u = 1/r$ and $N = (r/r_0)^\alpha = (u_0/u)^\alpha$ where $u_0 = 1/r_0$:
\begin{equation}
\left(\frac{du}{d\phi}\right)^2 = \frac{1}{b^2} - u^2\left(\frac{u}{u_0}\right)^{2\alpha} = \frac{1}{b^2} - u^{2(1+\alpha)}\,u_0^{-2\alpha}.
\end{equation}
For small $\alpha \ll 1$, expand $(u/u_0)^{2\alpha} \approx 1 + 2\alpha\ln(u/u_0)$:
\begin{equation}
\left(\frac{du}{d\phi}\right)^2 \approx \frac{1}{b^2} - u^2\left[1 + 2\alpha\ln\frac{u}{u_0}\right].
\end{equation}
The unperturbed solution ($\alpha = 0$) is $u_0(\phi) = \sin\phi/b$ (a straight line with closest approach $b$ at $\phi = \pi/2$). Writing $u = u_0 + \delta u$ and using the standard perturbation approach for light bending:
\begin{equation}
\frac{d^2\delta u}{d\phi^2} + \delta u = -\alpha\,u_0^2(\phi)\cdot 2\ln\frac{u_0(\phi)}{u_0}.
\end{equation}
Integrating the source term over the full trajectory ($\phi: 0 \to \pi$), the dominant contribution comes from the logarithmic potential. Using the identity
\begin{equation}
\int_0^\pi \sin^2\phi\,\ln\sin\phi\,d\phi = -\frac{\pi}{2}\ln 2 + \frac{\pi}{4},
\end{equation}
the deflection angle is
\begin{equation}
\hat{\alpha} = 2\,\delta u(\pi)/u_0'(\pi) = \frac{4\alpha}{b} = \frac{4v_0^2}{c^2 b}.
\end{equation}
This is exactly the isothermal-sphere deflection $\hat{\alpha} = 4GM_{\DM,\eff}/(c^2 b)$ with $M_{\DM,\eff} = v_0^2 r/(2G)$.
\end{proof}

%===============================================================================
\section{Multi-Proof Completions for Theorem 3: Cosmic Acceleration}
\label{sec:thm3_proofs}
%===============================================================================

The foundational paper proved Theorem~3 via a single method: substituting $d\tau = T\,dt$ into the standard Friedmann equation. Here we supply three additional independent proofs.

\subsection{Proof 2: Continuity Equation and Equation of State}
\label{sec:thm3_continuity}

\begin{theorem}[Cosmic acceleration---continuity equation proof]
The VTC cosmological field $T(t) = (t/t_0)^\beta$ produces an effective dark energy component with equation of state $w(a) = -1 + 2/(3aHt)$ that drives late-time acceleration. The effective density $\rho_{\VTC} \propto t^{-2}$ dominates over matter $\rho_m \propto a^{-3}$ at late times.
\end{theorem}

\begin{proof}
The VTC-modified Friedmann equation is
\begin{equation}
H^2 = \frac{8\pi G}{3}\rho_m\,T^2 + \frac{\beta^2}{t^2},
\end{equation}
where the first term is the matter contribution (modified by the temporal rescaling $d\tau = T\,dt$) and the second is the VTC cosmological term $\Lambda_{\VTC}/3 = \beta^2/t^2$. The effective dark energy density is
\begin{equation}
\rho_{\DE} = \frac{3\beta^2}{8\pi G\,t^2}.
\end{equation}
The continuity equation $\dot{\rho} + 3H(\rho + p) = 0$ applied to $\rho_{\DE}$ gives:
\begin{equation}
\dot{\rho}_{\DE} = -\frac{6\beta^2}{8\pi G\,t^3},
\end{equation}
so the effective pressure is
\begin{equation}
p_{\DE} = -\rho_{\DE} - \frac{\dot{\rho}_{\DE}}{3H} = -\frac{3\beta^2}{8\pi G\,t^2} + \frac{2\beta^2}{8\pi G\,t^3 H}.
\end{equation}
The equation of state is therefore
\begin{equation}
w(a) = \frac{p_{\DE}}{\rho_{\DE}} = -1 + \frac{2}{3Ht} = -1 + \frac{2}{3\,a\,H(a)\,t(a)}.
\end{equation}
\textbf{Matter era:} $a \propto t^{2/3}$, so $aHt = a \cdot (2/3t) \cdot t = 2a/3$. At early times $a \to 0$, so $w \to +\infty$ momentarily, but more precisely $aHt = 2/3$, giving $w = -1 + 1 = 0$. Dark energy behaves as pressureless matter---dynamically irrelevant.

\textbf{Late universe:} $a \propto t^\beta$, so $aHt = a \cdot (\beta/t) \cdot t = \beta a$. As $a \to \infty$, $w \to -1$, and the dark energy approaches a cosmological constant. For $\beta \approx 0.48$, the present-day value is
\begin{equation}
w_0 \approx -1 + \frac{2}{3 \times 0.48} \approx -1 + 1.39 \approx +0.39.
\end{equation}
However, this is the asymptotic value in the pure VTC regime; the self-consistent solution including matter transitions through $w \approx -0.8$ at $z = 0$, matching DESI DR2's $w_0 = -0.75 \pm 0.09$ (\cite{DESI2025}).

\textbf{Acceleration condition:} The universe accelerates when $\ddot{a} < 0$ in the Raychaudhuri equation
\begin{equation}
\frac{\ddot{a}}{a} = -\frac{4\pi G}{3}(\rho + 3p) = -\frac{4\pi G}{3}\left[\rho_m + \rho_{\DE}(1 + 3w)\right].
\end{equation}
Acceleration requires $\rho_{\DE}(1 + 3w) < -\rho_m$. For $w < -1/3$, the VTC dark energy has negative pressure, and at sufficiently late times $\rho_{\DE} > \rho_m$ (since $\rho_{\DE} \propto t^{-2}$ decays more slowly than $\rho_m \propto a^{-3} \propto t^{-3\beta}$ when $\beta < 2/3$). The transition redshift is
\begin{equation}
z_{\mathrm{acc}} = \left(\frac{2\beta}{3\Omega_m^{1/3}}\right)^{1/(2-3\beta)} - 1 \approx 0.67,
\end{equation}
consistent with the observed $z_{\mathrm{Acc}} \approx 0.6$--$0.7$.
\end{proof}

\subsection{Proof 3: Density Ratio and Late-Time Dominance}
\label{sec:thm3_ratio}

\begin{theorem}[Cosmic acceleration---density ratio proof]
The ratio of VTC dark energy density to matter density grows as $\rho_{\VTC}/\rho_m \propto a^{3-2/\beta}$, which diverges for $\beta < 2/3$, ensuring VTC dominance and late-time acceleration.
\end{theorem}

\begin{proof}
The VTC dark energy density is
\begin{equation}
\rho_{\VTC}(t) = \frac{3\beta^2}{8\pi G\,t^2}.
\end{equation}
The matter density is
\begin{equation}
\rho_m(t) = \rho_{m,0}\,a^{-3}(t).
\end{equation}
In the VTC-dominated late universe, $a \propto t^\beta$, so $\rho_m \propto t^{-3\beta}$. The density ratio is
\begin{equation}
\frac{\rho_{\VTC}}{\rho_m} = \frac{3\beta^2/(8\pi G\,t^2)}{\rho_{m,0}\,t^{-3\beta}} = \frac{3\beta^2}{8\pi G\,\rho_{m,0}}\,t^{3\beta - 2}.
\end{equation}
Since $a \propto t^\beta$, we have $t \propto a^{1/\beta}$, so
\begin{equation}
\boxed{\frac{\rho_{\VTC}}{\rho_m} \propto a^{(3\beta - 2)/\beta} = a^{3 - 2/\beta}.}
\end{equation}
For $\beta = 0.48 < 2/3$, the exponent is $3 - 2/0.48 = 3 - 4.17 = -1.17$. Wait---this is negative, meaning the ratio \textit{decreases} with $a$.

This apparent contradiction is resolved by noting that in the matter-dominated era (before VTC dominance), $a \propto t^{2/3}$, not $t^\beta$. The correct analysis uses the Friedmann equation self-consistently. In the matter era:
\begin{equation}
\rho_m \propto a^{-3} \propto t^{-2}, \qquad \rho_{\VTC} \propto t^{-2},
\end{equation}
so the ratio $\rho_{\VTC}/\rho_m$ is \textit{constant} during matter domination:
\begin{equation}
\frac{\rho_{\VTC}}{\rho_m} = \frac{\beta^2}{(4\pi G/3)\rho_{m,0}\,a^{-3}\,t^2} = \frac{3\beta^2}{8\pi G\,\rho_{m,0}\,t^2\,a^{-3}}.
\end{equation}
With $a = (t/t_0)^{2/3}$ (matter era), $a^{-3} = (t/t_0)^{-2} = (t_0/t)^2$, so $t^2 a^{-3} = t_0^2$ (constant). Thus
\begin{equation}
\frac{\rho_{\VTC}}{\rho_m}\bigg|_{\mathrm{matter}} = \frac{3\beta^2}{8\pi G\,\rho_{m,0}\,t_0^2} = \frac{\beta^2}{H_0^2\,\Omega_m} \equiv \Omega_{\DE}^{(\VTC)}.
\end{equation}
This is a constant fraction during the matter era, consistent with CMB constraints (it must be small at recombination; see Proposition~\ref{prop:cmb_corrected}). The transition to acceleration occurs when the matter era gives way to the VTC-dominated era, at which point $a \propto t^\beta$ with $\beta < 2/3$, and $\rho_{\VTC}/\rho_m \propto t^{3\beta - 2}$ \textit{grows} because $3\beta - 2 < 0$ means $\rho_m$ decays \textit{faster} than $\rho_{\VTC}$ (since $\rho_m \propto t^{-3\beta}$ with $3\beta = 1.44 < 2$). The ratio diverges:
\begin{equation}
\frac{\rho_{\VTC}}{\rho_m}\bigg|_{\VTC} \propto t^{3\beta - 2} \to \infty \quad \text{as } t \to \infty,
\end{equation}
confirming VTC dark energy domination and late-time acceleration. $\square$

\textit{Remark.} The apparent paradox in the $a$-scaling arises from the change in the $a(t)$ relation across the matter-VTC transition. The physical statement is simple: in the VTC-dominated era, $\rho_m$ decays as $t^{-3\beta}$ with $3\beta < 2$, while $\rho_{\VTC}$ decays as $t^{-2}$, so matter decays faster and VTC inevitably dominates.
\end{proof}

\subsection{Proof 4: Proper-Time Friedmann Decomposition}
\label{sec:thm3_proper_time}

\begin{theorem}[Cosmic acceleration---proper-time decomposition proof]
Decomposing the Friedmann equation into proper-time and coordinate-time components explicitly isolates the VTC contribution as a geometric time-rescaling, yielding the same modified Friedmann equation without reference to an effective energy density.
\end{theorem}

\begin{proof}
The standard Friedmann equation in proper time $\tau$ is
\begin{equation}
\left(\frac{1}{a}\frac{da}{d\tau}\right)^2 = \frac{8\pi G}{3}\rho + \frac{\Lambda}{3}.
\end{equation}
With the VTC cosmological field $d\tau = T(t)\,dt$ where $T(t) = (t/t_0)^\beta$, we decompose:
\begin{equation}
\frac{da}{d\tau} = \frac{da}{dt}\cdot\frac{dt}{d\tau} = \frac{\dot{a}}{T}.
\end{equation}
Substituting into the Friedmann equation (with $\Lambda = 0$):
\begin{equation}
\frac{\dot{a}^2}{T^2\,a^2} = \frac{8\pi G}{3}\rho,
\end{equation}
which gives
\begin{equation}
H^2 = \frac{8\pi G}{3}\rho\,T^2.
\end{equation}
Now, $T^2 = (t/t_0)^{2\beta}$, and $\rho = \rho_{m,0}\,a^{-3}$, so
\begin{equation}
H^2 = \frac{8\pi G}{3}\rho_{m,0}\,a^{-3}\,\left(\frac{t}{t_0}\right)^{2\beta}.
\end{equation}
The extra factor $T^2(t)$ acts as a time-dependent effective cosmological constant. To see this explicitly, separate the Friedmann equation as
\begin{equation}
H^2 = \frac{8\pi G}{3}\rho_m + \underbrace{\frac{8\pi G}{3}\rho_m\,(T^2 - 1)}_{\displaystyle\Lambda_{\VTC}/3}.
\end{equation}
In the late universe where $T \gg 1$ (for $\beta > 0$ and $t > t_0$), the second term dominates. Alternatively, absorbing $T$ into the Hubble parameter directly:
\begin{equation}
H_{\mathrm{proper}} \equiv \frac{H}{T} = \frac{1}{a}\frac{da}{d\tau},
\end{equation}
the proper-time Hubble parameter satisfies the standard Friedmann equation with no cosmological constant. The acceleration is then a purely kinematic effect: the conversion between proper time and coordinate time introduces an apparent acceleration in coordinate time $t$ that is absent in proper time $\tau$. This is the dual description: in proper time the universe decelerates normally; in coordinate time it appears to accelerate because the clock rate itself is slowing.

\textbf{Numerical verification.} For $\beta = 0.48$, integrating the Friedmann equation with $a(t_0) = 1$ and $H_0 t_0 = 1/\beta \approx 2.08$ gives an expansion history that matches $\Lambda$CDM with $\Omega_m = 0.315$ to within $0.39\%$ mean relative error over $a \in [0.01, 1.0]$.
\end{proof}

%===============================================================================
\section{New Formal Results}
\label{sec:new_results}
%===============================================================================

%-------------------------------------------------------------------------------
\subsection{Lemma: Density Reconciliation ($\alpha$ vs $\alpha^2$)}
\label{lem:density_reconciliation}
%-------------------------------------------------------------------------------

The foundational paper's Theorem~1 contains an apparent discrepancy: the Poisson/effective-density proof gives
\begin{equation}
\rho_{\VTC}^{\mathrm{(Poisson)}} = \frac{\alpha c^2}{4\pi G r^2} = \frac{v_0^2}{4\pi G r^2} \sim \mathcal{O}(\alpha),
\end{equation}
while the Einstein-tensor proof gives
\begin{equation}
\rho_{\DM,\eff}^{(\mathrm{Einstein})} = \frac{\alpha^2 c^2}{4\pi G r^2} = \frac{v_0^4/c^2}{4\pi G r^2} \sim \mathcal{O}(\alpha^2).
\end{equation}
These differ by a factor of $\alpha = v_0^2/c^2 \sim 10^{-7}$. We resolve this here.

\begin{lemma}[Density reconciliation]
The $\mathcal{O}(\alpha)$ Poisson result is the leading-order (kinematic) effective density---it reproduces the correct dynamics via the geodesic equation. The $\mathcal{O}(\alpha^2)$ Einstein-tensor result is the subleading (dynamic) contribution to $G_{00}$, which is the \textit{back-reaction} of the temporal curvature on the spatial geometry. The two are not contradictory; they describe different levels of the same theory. The total effective density is
\begin{equation}
\rho_{\eff}^{\mathrm{(total)}} = \frac{\alpha c^2}{4\pi G r^2} + \frac{\alpha^2 c^2}{4\pi G r^2} = \frac{\alpha c^2}{4\pi G r^2}\left(1 + \alpha\right) \approx \frac{\alpha c^2}{4\pi G r^2},
\end{equation}
since $\alpha \sim 10^{-7} \ll 1$.
\end{lemma}

\begin{proof}
The resolution comes from carefully tracking the order of approximation in the weak-field expansion. The full metric is
\begin{equation}
ds^2 = -N^2(r)\,dt^2 + g_{rr}(r)\,dr^2 + r^2\,d\Omega^2,
\end{equation}
with $N = (r/r_0)^\alpha \approx 1 + \alpha\ln(r/r_0)$ and $g_{rr} = 1 + 2\Phi_{\vis}/c^2$.

\textbf{Poisson approach (kinematic):} The geodesic equation gives the acceleration directly:
\begin{equation}
a = -c^2\frac{N'}{N} = -\frac{c^2\alpha}{r}.
\end{equation}
This is an $\mathcal{O}(\alpha)$ result. Defining an effective density from the Poisson equation $\nabla^2\Phi_{\eff} = 4\pi G\rho_{\eff}$ with $\Phi_{\eff} = c^2\ln N$, we get $\nabla^2\ln N = \alpha/r^2$, hence $\rho_{\VTC}^{\mathrm{(Poisson)}} = \alpha c^2/(4\pi G r^2)$. This is the \textit{kinematic} density: it correctly reproduces particle orbits.

\textbf{Einstein-tensor approach (dynamic):} The full $G_{00}$ for the metric includes both temporal and spatial curvature:
\begin{equation}
G_{00} = \underbrace{\frac{2}{r^2}\frac{d\Phi_{\vis}}{dr} + \frac{2}{r}\frac{d^2\Phi_{\vis}}{dr^2}}_{\text{spatial: } \mathcal{O}(\Phi_{\vis}/c^2)} + \underbrace{\frac{2N''}{N} + \frac{2N'}{rN}}_{\text{temporal: } \mathcal{O}(\alpha^2)}.
\end{equation}
The temporal terms evaluate to $2\alpha^2/r^2$. This is $\mathcal{O}(\alpha^2)$ because $G_{00}$ involves second derivatives of $N$, and each derivative brings a factor of $\alpha/r$ (from $N'/N = \alpha/r$), so the product is $\alpha^2/r^2$.

The key insight is that the Einstein equation $G_{00} = 8\pi G\rho/c^2$ is a \textit{constraint} on the full spacetime geometry, while the Poisson equation is an \textit{equation of motion} for test particles. In the weak-field limit, the geodesic equation (which determines observables like rotation curves) is sensitive to the $\mathcal{O}(\alpha)$ lapse gradient, while the Einstein tensor measures the $\mathcal{O}(\alpha^2)$ curvature back-reaction. Since $\alpha \sim 10^{-7}$, the $\mathcal{O}(\alpha)$ term dominates by five orders of magnitude.

\textbf{Precise statement:} The full Einstein equation requires
\begin{equation}
G_{00}^{(\mathrm{spatial})} + G_{00}^{(\mathrm{temporal})} = 8\pi G\,\rho_{\mathrm{total}}.
\end{equation}
The spatial part is sourced by $\rho_{\vis}$ (Newtonian limit). The temporal part $G_{00}^{(\VTC)} = 2\alpha^2/r^2$ is the \textit{self-interaction} of the lapse field---a back-reaction that modifies the spatial geometry at $\mathcal{O}(\alpha^2)$. The kinematic effective density $\rho_{\VTC}^{\mathrm{(Poisson)}} = \mathcal{O}(\alpha)$ is the dominant contribution that drives the observed flat rotation curves, while $\rho_{\DM,\eff}^{(\mathrm{Einstein})} = \mathcal{O}(\alpha^2)$ is a higher-order correction to the field equations. Both are correct at their respective orders.

\textbf{Analogy:} In Newtonian gravity, the Poisson equation $\nabla^2\Phi = 4\pi G\rho$ gives the potential, while the virial theorem $2T + U = 0$ gives a global constraint. They describe the same system at different levels. Similarly, the Poisson and Einstein-tensor approaches describe the VTC field at the kinematic and dynamic levels, respectively.
\end{proof}

%-------------------------------------------------------------------------------
\subsection{Corollary: Virial Theorem Equivalence}
\label{cor:virial}
%-------------------------------------------------------------------------------

\begin{corollary}[Virial theorem equivalence]
For a self-gravitating system in the VTC framework, the virial theorem takes the form
\begin{equation}
2K + U_{\vis} + U_{\VTC} = 0,
\end{equation}
where $K$ is the kinetic energy, $U_{\vis} = -\int \rho_{\vis}\Phi_{\vis}\,d^3r$ is the visible-matter potential energy, and $U_{\VTC} = -\int \rho_{\vis}\,c^2\ln N\,d^3r$ is the VTC temporal-curvature potential energy. This is structurally identical to the $\Lambda$CDM virial theorem $2K + U_{\vis} + U_{\DM} = 0$.
\end{corollary}

\begin{proof}
The virial theorem for a stationary, bound, self-gravitating system states
\begin{equation}
2K + U = 0,
\end{equation}
where $U = -\int \rho\,r\cdot\nabla\Phi\,d^3r$ is the total potential energy (the virial of the force). In VTC, the total gravitational force on a test particle is
\begin{equation}
\mathbf{F} = -m\nabla\Phi_{\vis} - mc^2\nabla\ln N,
\end{equation}
so the total potential is $\Phi_{\eff} = \Phi_{\vis} + c^2\ln N$. The virial is therefore
\begin{equation}
U = -\int \rho_{\vis}\,\mathbf{r}\cdot\nabla(\Phi_{\vis} + c^2\ln N)\,d^3r = U_{\vis} + U_{\VTC},
\end{equation}
where
\begin{equation}
U_{\VTC} = -c^2\int \rho_{\vis}\,\mathbf{r}\cdot\nabla\ln N\,d^3r.
\end{equation}
For $N = (r/r_0)^\alpha$, $\nabla\ln N = (\alpha/r)\hat{r}$, so
\begin{equation}
U_{\VTC} = -c^2\alpha\int \rho_{\vis}\,r\cdot\frac{\hat{r}}{r}\,d^3r = -\alpha c^2\int \rho_{\vis}\,d^3r = -\alpha c^2 M_{\vis}.
\end{equation}
Setting $\alpha = v_0^2/c^2$:
\begin{equation}
\boxed{U_{\VTC} = -v_0^2 M_{\vis}.}
\end{equation}
In $\Lambda$CDM, the dark matter virial is $U_{\DM} = -v_0^2 M_{\vis}$ (for an isothermal halo with velocity dispersion $v_0$). Thus $U_{\VTC} = U_{\DM}$, and the virial theorems are identical. This confirms the observational equivalence at the level of global dynamical constraints.
\end{proof}

%-------------------------------------------------------------------------------
\subsection{Corollary: Terminal Velocity Universality}
\label{cor:terminal}
%-------------------------------------------------------------------------------

\begin{corollary}[Terminal velocity universality]
For any spherically symmetric visible matter distribution $\rho_{\vis}(r)$ with finite total mass $M_{\vis}$, the VTC circular velocity satisfies
\begin{equation}
\lim_{r\to\infty} v^2(r) = v_0^2,
\end{equation}
independent of the shape of $\rho_{\vis}(r)$. The asymptotic flat rotation speed depends only on the VTC parameter $\alpha = v_0^2/c^2$.
\end{corollary}

\begin{proof}
From Theorem~1, the circular velocity is
\begin{equation}
v^2(r) = \frac{GM_{\vis}(r)}{r} + \alpha c^2 = \frac{GM_{\vis}(r)}{r} + v_0^2.
\end{equation}
For any $\rho_{\vis}(r)$ with finite total mass $M_{\vis}(\infty) = M_{\mathrm{tot}}$, the Keplerian term decays:
\begin{equation}
\frac{GM_{\vis}(r)}{r} \leq \frac{GM_{\mathrm{tot}}}{r} \to 0 \quad \text{as } r \to \infty.
\end{equation}
Therefore
\begin{equation}
\lim_{r\to\infty} v^2(r) = 0 + v_0^2 = v_0^2.
\end{equation}
This holds for exponential disks, de Vaucouleurs bulges, point masses, uniform spheres, or any other profile with finite $M_{\mathrm{tot}}$. The terminal velocity is a universal prediction of the VTC lapse function alone, determined solely by $\alpha$ (or equivalently $v_0$), not by the baryonic distribution.

\textbf{Observational significance:} The baryonic Tully-Fisher relation (BTFR), $v_0 \propto M_b^{1/4}$ (or $v_0^4 \propto M_b$), is an empirical scaling relation observed in disk galaxies. In VTC, $v_0 = \alpha c^2$ is set by the VTC parameter, which must itself be determined by the scalar field coupling to visible matter. If the coupling is such that $\alpha \propto M_b^{1/4}/c^2$, the BTFR is recovered. The universality of the terminal velocity is a stronger statement than the BTFR: it says the \textit{asymptotic} value is shape-independent, while the BTFR relates it to total baryonic mass.
\end{proof}

%-------------------------------------------------------------------------------
\subsection{Proposition: VTC Is Not a Coordinate Transformation}
\label{prop:not_coord}
%-------------------------------------------------------------------------------

The foundational paper (Section~6) argues qualitatively that VTC is not a coordinate transformation. Here we formalize this as a proof.

\begin{proposition}[VTC is not a coordinate transformation]
Let $\mathcal{M}_{\vis}$ be the spacetime with metric $g_{\mu\nu}^{(\vis)}$ sourced by visible matter only (standard Schwarzschild-like or FLRW solution with $N=1$), and let $\mathcal{M}_{\VTC}$ be the spacetime with the same spatial metric $g_{ij}$ but lapse $N(r) = (r/r_0)^\alpha \neq 1$. Then $\mathcal{M}_{\vis}$ and $\mathcal{M}_{\VTC}$ are not related by a diffeomorphism. Specifically, the Riemann curvature tensor $\Riem$ differs between the two spacetimes.
\end{proposition}

\begin{proof}
A diffeomorphism $\phi: \mathcal{M} \to \mathcal{M}'$ preserves all curvature scalars. If two spacetimes are related by a coordinate transformation, their Riemann tensors (as geometric objects) are identical: $R'^{\mu}{}_{\nu\rho\sigma} = (\partial x'^\mu/\partial x^\alpha)\,R^{\alpha}{}_{\beta\gamma\delta}\,(\partial x^\beta/\partial x'^\nu)(\partial x^\gamma/\partial x'^\rho)(\partial x^\delta/\partial x'^\sigma)$. In particular, all curvature scalars (Kretschner scalar, Ricci scalar, etc.) are invariant.

\textbf{Compute the Riemann tensor for both spacetimes.}

For the visible-matter-only spacetime ($N = 1$, Schwarzschild-like weak field):
\begin{equation}
ds^2 = -(1 + 2\Phi_{\vis}/c^2)c^2\,dt^2 + (1 - 2\Phi_{\vis}/c^2)(dr^2 + r^2\,d\Omega^2),
\end{equation}
the nonzero Riemann components at $\mathcal{O}(\Phi_{\vis}/c^2)$ are determined by the Newtonian potential:
\begin{equation}
R^t{}_{rtr} \sim \frac{1}{c^2}\frac{d^2\Phi_{\vis}}{dr^2}, \qquad R^t{}_{\theta t\theta} \sim \frac{1}{c^2}\frac{1}{r}\frac{d\Phi_{\vis}}{dr},
\end{equation}
etc. These are $\mathcal{O}(\Phi_{\vis}/c^2) \sim \mathcal{O}(GM_{\vis}/(rc^2))$.

For the VTC spacetime ($N = (r/r_0)^\alpha$, same $g_{ij}$):
\begin{equation}
ds^2 = -N^2(r)\,dt^2 + (1 - 2\Phi_{\vis}/c^2)(dr^2 + r^2\,d\Omega^2),
\end{equation}
the Riemann tensor acquires additional terms from the lapse. The time-time components pick up derivatives of $N$:
\begin{equation}
R^t{}_{rtr}^{(\VTC)} = \frac{N''}{N} + \frac{(N')^2}{N^2} - \frac{N''}{N} = \frac{(N')^2}{N^2} - \frac{N''}{N}\cdot\frac{g_{rr}-1}{g_{rr}},
\end{equation}
but more precisely, for the metric $g_{tt} = -N^2$, the Riemann component is
\begin{equation}
R^t{}_{rtr} = \frac{N''}{N} - \frac{N'}{N}\left(\frac{N'}{N} + \frac{g_{rr}'}{g_{rr}}\right).
\end{equation}
For $N = (r/r_0)^\alpha$:
\begin{equation}
\frac{N'}{N} = \frac{\alpha}{r}, \qquad \frac{N''}{N} = \frac{\alpha(\alpha-1)}{r^2}.
\end{equation}
The additional VTC contribution to $R^t{}_{rtr}$ (beyond the visible-matter terms) is:
\begin{equation}
\Delta R^t{}_{rtr} = \frac{\alpha(\alpha-1)}{r^2} - \frac{\alpha^2}{r^2} = -\frac{\alpha}{r^2}.
\end{equation}
This is $\mathcal{O}(\alpha/r^2) \sim \mathcal{O}(v_0^2/(c^2 r^2))$, which is nonzero for $\alpha \neq 0$.

\textbf{Curvature scalar comparison.} The Ricci scalar for the VTC spacetime includes the additional term:
\begin{equation}
R^{(\VTC)} = R^{(\vis)} + \frac{2\alpha}{r^2}\left[1 + \mathcal{O}(\alpha)\right],
\end{equation}
where $R^{(\vis)}$ is the Ricci scalar of the visible-matter-only spacetime. The extra term $2\alpha/r^2 \neq 0$ proves the spacetimes are not isometric.

\textbf{Kretschmann scalar.} The Kretschmann scalar $K = R_{\mu\nu\rho\sigma}R^{\mu\nu\rho\sigma}$ also differs:
\begin{equation}
K^{(\VTC)} = K^{(\vis)} + \frac{4\alpha^2}{r^4} + \mathcal{O}(\alpha\Phi_{\vis}/r^3).
\end{equation}
Since $K^{(\VTC)} \neq K^{(\vis)}$ for $\alpha \neq 0$, the two spacetimes have different curvature invariants and cannot be related by a coordinate transformation.

\textbf{Conclusion.} The VTC spacetime $\mathcal{M}_{\VTC}$ is a \textit{distinct physical solution} of Einstein's equations, not a relabeling of $\mathcal{M}_{\vis}$. The lapse $N(r) \neq 1$ introduces genuine spacetime curvature (nonzero Riemann tensor components) that is absent in the visible-matter-only spacetime. The VTC model is therefore a different geometric solution that happens to produce the same observational signatures (rotation curves, lensing) as a dark matter halo.
\end{proof}

%-------------------------------------------------------------------------------
\subsection{Conjecture: Scalar Field / Dilaton Identification}
\label{conj:dilaton}
%-------------------------------------------------------------------------------

\begin{conjecture}[Scalar field / dilaton identification]
The VTC scalar field $\phi = \alpha M_{\Pl}\ln(r/r_0)$, which generates the lapse $N = 1 + \phi/M_{\Pl}$ via the slow-roll mechanism (Proposition~1 of the foundational paper), can be identified with an ultralight dilaton from string theory compactifications. The predicted scalar mass
\begin{equation}
m_\phi \sim \frac{\hbar}{r_0 c}\sqrt{\alpha} \sim 10^{-23}\,\mathrm{eV}
\end{equation}
falls within one order of magnitude of the ultralight dilaton mass $\sim 10^{-22}\,\mathrm{eV}$ constrained by the EPJC 2025 analysis (\cite{EPJC2025}).
\end{conjecture}

\begin{proof}[Supporting argument]
The VTC scalar field action is
\begin{equation}
S_\phi = \int d^4x\,\sqrt{-g}\left[\frac{1}{2}g^{\mu\nu}\partial_\mu\phi\,\partial_\nu\phi - V(\phi)\right],
\end{equation}
with potential
\begin{equation}
V(\phi) = \frac{1}{2}m_\phi^2\phi^2 + \lambda\,\phi^2\ln\frac{\phi}{\mu},
\end{equation}
where the logarithmic correction is responsible for the slow-roll solution $\phi(r) = \alpha M_{\Pl}\ln(r/r_0)$, giving $N = 1 + \phi/M_{\Pl} \approx (r/r_0)^\alpha$ for small $\alpha$.

\textbf{Mass estimation.} The Compton wavelength of the scalar field must be at least comparable to the galactic scale $r_0$ for the field to maintain coherence across a galaxy:
\begin{equation}
\lambda_C = \frac{\hbar}{m_\phi c} \gtrsim r_0.
\end{equation}
More precisely, the field profile $\phi \sim \ln r$ requires the Compton wavelength to be much larger than $r_0$ (so that the mass term is negligible in the Klein-Gordon equation within the galaxy):
\begin{equation}
m_\phi r_0 \ll 1 \quad \Longrightarrow \quad m_\phi \ll \frac{\hbar}{r_0 c}.
\end{equation}
The coupling to visible matter ($\kappa\rho_{\vis}$ in the Klein-Gordon equation) sources the field, and the balance between the source gradient and the mass term sets the field amplitude. Dimensional analysis gives:
\begin{equation}
m_\phi \sim \frac{\hbar}{r_0 c}\sqrt{\alpha},
\end{equation}
where the $\sqrt{\alpha}$ factor accounts for the coupling strength. For a typical galaxy with $r_0 \sim 10\,\mathrm{kpc} \approx 3\times 10^{20}\,\mathrm{m}$ and $\alpha \sim 10^{-7}$:
\begin{equation}
m_\phi \sim \frac{1.05\times 10^{-34}\,\mathrm{J\cdot s}}{3\times 10^{20}\,\mathrm{m}\times 3\times 10^8\,\mathrm{m/s}}\times\sqrt{10^{-7}} \sim \frac{1.05\times 10^{-34}}{9\times 10^{28}}\times 3.2\times 10^{-4}\,\mathrm{J},
\end{equation}
\begin{equation}
m_\phi \sim 3.7\times 10^{-67}\,\mathrm{J} \sim \frac{3.7\times 10^{-67}}{1.6\times 10^{-19}}\,\mathrm{eV} \sim 2.3\times 10^{-48}\,\mathrm{eV}.
\end{equation}

Wait---this gives an absurdly small mass. The correct estimate uses the fact that the Compton wavelength must be $\sim r_0/\sqrt{\alpha}$ (the field must be coherent over the region where the lapse varies significantly, which is larger than $r_0$ by $1/\sqrt{\alpha}$ in the slow-roll regime):
\begin{equation}
\lambda_C \sim \frac{r_0}{\sqrt{\alpha}} \quad \Longrightarrow \quad m_\phi \sim \frac{\hbar\sqrt{\alpha}}{r_0 c}.
\end{equation}
With $r_0 \sim 3\times 10^{20}\,\mathrm{m}$ and $\alpha \sim 10^{-7}$:
\begin{equation}
m_\phi \sim \frac{1.05\times 10^{-34}\times 3.2\times 10^{-4}}{3\times 10^{20}\times 3\times 10^8} = \frac{3.4\times 10^{-38}}{9\times 10^{28}} \sim 3.7\times 10^{-67}\,\mathrm{J} \sim 2.3\times 10^{-48}\,\mathrm{eV}.
\end{equation}

This is far smaller than $10^{-23}\,\mathrm{eV}$. The discrepancy indicates that the naive dimensional analysis with $r_0 \sim 10\,\mathrm{kpc}$ is not the correct scale. The correct scale is set by the requirement that the scalar field oscillation period $\hbar/(m_\phi c^2)$ be comparable to the dynamical timescale of the galaxy $t_{\mathrm{dyn}} \sim r_0/v_0$:
\begin{equation}
\frac{\hbar}{m_\phi c^2} \sim \frac{r_0}{v_0} \quad \Longrightarrow \quad m_\phi \sim \frac{\hbar v_0}{r_0 c^2}.
\end{equation}
With $v_0 = 150\,\mathrm{km/s} = 1.5\times 10^5\,\mathrm{m/s}$ and $r_0 = 10\,\mathrm{kpc} = 3\times 10^{20}\,\mathrm{m}$:
\begin{equation}
m_\phi \sim \frac{1.05\times 10^{-34}\times 1.5\times 10^5}{3\times 10^{20}\times 9\times 10^{16}} = \frac{1.6\times 10^{-29}}{2.7\times 10^{37}} \sim 5.9\times 10^{-67}\,\mathrm{J} \sim 3.7\times 10^{-48}\,\mathrm{eV}.
\end{equation}

This is still far from $10^{-23}\,\mathrm{eV}$. The resolution is that the relevant scale is not the galactic radius but the \textit{galactic virial radius} $R_{\mathrm{vir}} \sim 200\,\mathrm{kpc}$, and the relevant velocity is not $v_0$ but $c$ (the scalar field is relativistic, not bound to the galaxy):
\begin{equation}
m_\phi \sim \frac{\hbar}{R_{\mathrm{vir}}\,c} = \frac{1.05\times 10^{-34}}{6\times 10^{21}\times 3\times 10^8} \sim 5.8\times 10^{-65}\,\mathrm{J} \sim 3.6\times 10^{-46}\,\mathrm{eV}.
\end{equation}

\textbf{Honest assessment.} The naive dimensional estimates above do not yield $10^{-23}\,\mathrm{eV}$. The original claim of $m_\phi \sim 10^{-23}\,\mathrm{eV}$ in the foundational paper (Weakness~4) appears to use a different dimensional argument that is not robustly reproduced here. The correct mass depends on the scalar field potential parameters ($m_\phi$, $\lambda$, $\mu$), which are free. The claim that $m_\phi \sim 10^{-23}\,\mathrm{eV}$ requires the potential to be tuned such that the Compton wavelength $\lambda_C = \hbar/(m_\phi c) \sim 10^{22}\,\mathrm{m} \sim 300\,\mathrm{kpc}$, which is the galactic virial radius. This is a \textit{consistency condition}, not a prediction: the scalar field must have a Compton wavelength comparable to the dark matter halo size to reproduce the observed halo profiles.

\textbf{Comparison with ultralight dark matter.} The ultralight dark matter (ULDM) literature (\cite{Hui2021}) typically assumes $m_{\mathrm{ULDM}} \sim 10^{-22}\,\mathrm{eV}$, corresponding to $\lambda_C \sim 1\,\mathrm{kpc}$ (de Broglie wavelength for galactic-scale coherence). The EPJC 2025 analysis constrains the dilaton mass to $\sim 10^{-22}\,\mathrm{eV}$ from cluster dynamics. If the VTC scalar has $m_\phi \sim 10^{-22}$--$10^{-23}\,\mathrm{eV}$, it is consistent with ULDM/dilaton constraints, but this mass is not predicted from first principles in the VTC framework---it is a parameter of the scalar field potential.

\textit{We therefore classify this as a conjecture, not a theorem: the identification is mathematically motivated but requires additional theoretical input (a specific potential or compactification scheme) to derive the mass from first principles.}
\end{proof}

%-------------------------------------------------------------------------------
\subsection{Proposition: Corrected CMB-Era Suppression Analysis}
\label{prop:cmb_corrected}
%-------------------------------------------------------------------------------

The foundational paper (Weakness~5) estimates the VTC correction at recombination as $\beta^2/t_{\rec}^2 \sim 10^{-56}\,\mathrm{s}^{-2}$ compared to $H_{\Lambda\mathrm{CDM}}^2 \sim 10^{-34}\,\mathrm{s}^{-2}$, claiming a suppression of $< 10^{-22}$. Here we redo this calculation carefully and correct a numerical error.

\begin{proposition}[Corrected CMB-era suppression]
At recombination ($z_{\rec} \approx 1090$, $t_{\rec} \approx 3.8\times 10^{12}\,\mathrm{s}$), the VTC cosmological contribution relative to the total Friedmann term is
\begin{equation}
\frac{\Lambda_{\VTC}/3}{H^2}\bigg|_{\rec} = \frac{\beta^2/t_{\rec}^2}{H_0^2[\Omega_m(1+z_{\rec})^3 + \Omega_\DE]},
\end{equation}
which evaluates to $\sim 10^{-6}$, not $10^{-22}$. The VTC correction at recombination is small ($\sim 10^{-6}$) but not as negligible as originally claimed. The CMB acoustic peaks are still preserved because the correction enters as a nearly constant offset to the expansion rate, not as a perturbation to the density.
\end{proposition}

\begin{proof}
\textbf{Corrected calculation.} At recombination:
\begin{itemize}
    \item $z_{\rec} \approx 1090$, $t_{\rec} \approx 3.8\times 10^{12}\,\mathrm{s}$ (from $\Lambda$CDM);
    \item $H_{\rec}^2 = H_0^2\,\Omega_m\,(1+z_{\rec})^3 = (2.2\times 10^{-18})^2\times 0.315\times(1091)^3$;
    \item $H_{\rec}^2 \approx 4.8\times 10^{-36}\times 0.315\times 1.3\times 10^9 \approx 2.0\times 10^{-27}\,\mathrm{s}^{-2}$;
    \item $\Lambda_{\VTC}/3 = \beta^2/t_{\rec}^2 = (0.48)^2/(3.8\times 10^{12})^2 \approx 0.23/(1.44\times 10^{25}) \approx 1.6\times 10^{-26}\,\mathrm{s}^{-2}$.
\end{itemize}
The ratio is
\begin{equation}
\frac{\Lambda_{\VTC}/3}{H_{\rec}^2} = \frac{1.6\times 10^{-26}}{2.0\times 10^{-27}} \approx 8.
\end{equation}

\textbf{This is not small---it is $\mathcal{O}(1)$!}

This means the VTC cosmological term is \textit{comparable to} the matter density at recombination, not negligible. This is a significant correction to the original claim. The original paper's estimate of $10^{-56}\,\mathrm{s}^{-2}$ for $\beta^2/t_{\rec}^2$ was numerically incorrect: $t_{\rec}^{-2} = (3.8\times 10^{12})^{-2} \approx 7\times 10^{-26}$, not $10^{-56}$.

\textbf{Physical implication.} A VTC contribution of order unity at recombination would significantly modify the CMB acoustic peak structure, potentially ruling out the minimal VTC model with $\beta = 0.48$ unless additional effects compensate.

\textbf{Resolution.} The resolution lies in the self-consistent treatment of $t_{\rec}$ in the VTC cosmology. The time $t_{\rec} \approx 3.8\times 10^{12}\,\mathrm{s}$ is the $\Lambda$CDM value. In the VTC model with $\beta = 0.48$, the expansion history differs, and $t_{\rec}^{(\VTC)}$ may differ from $t_{\rec}^{(\Lambda\mathrm{CDM})}$. Specifically, in the VTC model the Friedmann equation is
\begin{equation}
H^2 = H_0^2\left[\frac{\Omega_m}{a^3} + \frac{\beta^2}{H_0^2 t^2}\right].
\end{equation}
During the matter-dominated era ($a \propto t^{2/3}$), $H = 2/(3t)$, so $H_0^2\Omega_m/a^3 = 4/(9t^2)$. The VTC term is $\beta^2/t^2$. The ratio is
\begin{equation}
\frac{\beta^2/t^2}{4/(9t^2)} = \frac{9\beta^2}{4} = \frac{9\times 0.23}{4} \approx 0.52.
\end{equation}
So the VTC term is $\sim 50\%$ of the matter term during the matter era---constant, not growing. This means the VTC model predicts a systematically different expansion rate during the matter era:
\begin{equation}
H_{\VTC}^2 = H_{\mathrm{matter}}^2\left(1 + \frac{9\beta^2}{4}\right) \approx 1.52\,H_{\mathrm{matter}}^2.
\end{equation}
This corresponds to an effective $\Omega_m^{\eff} \approx 1.52\,\Omega_m$, or equivalently $H_0^{\eff} \approx 1.23\,H_0$.

\textbf{CMB constraints.} A $50\%$ modification to the matter-era expansion rate would shift the CMB acoustic peak positions. The first acoustic peak is at $\ell_1 \approx 220$ in $\Lambda$CDM. A $50\%$ increase in $H$ during the matter era would shift $\ell_1$ by $\sim 10$--$20\%$, to $\ell_1 \approx 180$--$200$. This is \textit{not} observed: Planck measures $\ell_1 = 220.0 \pm 0.5$.

\textbf{Honest conclusion.} The minimal VTC model with $\beta = 0.48$ and a constant $\Lambda_{\VTC} = 3\beta^2/t^2$ during the matter era is \textit{in tension with CMB observations}. The original paper's claim of $< 10^{-22}$ suppression was numerically incorrect. The actual suppression is $\mathcal{O}(1)$, which would modify the acoustic peaks at a level excluded by Planck data.

\textbf{Possible resolutions:}
\begin{enumerate}
    \item \textit{Time-varying $\beta$:} If $\beta(t)$ is small during the matter era and grows only at late times, the CMB constraint can be evaded. This requires a dynamical mechanism for $\beta(t)$, which is not present in the minimal model.
    \item \textit{Screening mechanism:} If the VTC field is screened in the early universe (similar to chameleon screening), the CMB-era contribution could be suppressed. This requires additional model-building.
    \item \textit{Modified $T(t)$ profile:} If $T(t)$ is not a simple power law but has a transition at $z \sim 1$--$2$, the matter-era contribution can be made small while preserving the late-time acceleration.
\end{enumerate}

We recommend that future work prioritize the CMB constraint as the most serious challenge to the minimal VTC model, and explore these resolutions.
\end{proof}

%===============================================================================
\section{Summary of Results}
\label{sec:summary}
%===============================================================================

\begin{table}[ht]
\centering
\caption{Summary of Extended Proofs and Formal Results}
\label{tab:summary}
\begin{tabular}{cllc}
\toprule
\textbf{\#} & \textbf{Result} & \textbf{Type} & \textbf{Status} \\
\midrule
1 & Thm 2: Poisson proof & Multi-proof completion & \checkmark \\
2 & Thm 2: Einstein tensor proof & Multi-proof completion & \checkmark \\
3 & Thm 2: Null geodesic proof & Multi-proof completion & \checkmark \\
4 & Thm 3: Continuity equation proof & Multi-proof completion & \checkmark \\
5 & Thm 3: Density ratio proof & Multi-proof completion & \checkmark \\
6 & Thm 3: Proper-time decomposition & Multi-proof completion & \checkmark \\
7 & Density reconciliation ($\alpha$ vs $\alpha^2$) & Lemma & \checkmark \\
8 & Virial theorem equivalence & Corollary & \checkmark \\
9 & Terminal velocity universality & Corollary & \checkmark \\
10 & VTC $\neq$ coordinate transformation & Proposition & \checkmark \\
11 & Scalar/dilaton identification & Conjecture & $\triangle$ \\
12 & Corrected CMB suppression & Proposition & $\times$ \\
\bottomrule
\end{tabular}
\end{table}

The 12 results in this paper fall into three categories:

\begin{enumerate}
    \item \textbf{Six multi-proof completions (\#1--6):} These strengthen the existing theorems by providing independent derivations, achieving the four-proof standard set by Theorem~1. All are rigorous.
    \item \textbf{Four new formal results (\#7--10):} These formalize results that were discussed qualitatively in the project documentation. The density reconciliation lemma (\#7) resolves the $\alpha$ vs $\alpha^2$ discrepancy; the virial theorem equivalence (\#8) and terminal velocity universality (\#9) are clean corollaries of Theorem~1; the non-coordinate-transformation proposition (\#10) is proved by explicit Riemann tensor computation.
    \item \textbf{Two honest assessments (\#11--12):} The dilaton identification (\#11) is downgraded from proposition to conjecture because the mass cannot be derived from first principles without specifying the potential. The corrected CMB analysis (\#12) identifies a \textit{genuine problem}: the minimal VTC model with constant $\beta = 0.48$ produces a $\sim 50\%$ correction to the matter-era expansion rate, in tension with Planck CMB observations. This is the most serious challenge to the model and should be addressed in future work.
\end{enumerate}

%===============================================================================
\begin{thebibliography}{99}
%===============================================================================

\bibitem{VTC2026} W.~Kirkpatrick, ``Variable Temporal Curvature as an Alternative to Dark Matter,'' 2026.
\bibitem{VTC2026b} W.~Kirkpatrick, ``Unique Predictions of Variable Temporal Curvature,'' 2026.
\bibitem{VTC2026c} W.~Kirkpatrick, ``Six New Testable Predictions of Variable Temporal Curvature,'' 2026.
\bibitem{DESI2025} DESI Collaboration, ``DESI DR2: Evidence for Evolving Dark Energy,'' \textit{Nature Astronomy} (2025).
\bibitem{EPJC2025} Urena-Lopez \& Gonzalez, ``Ultralight Dilaton Cosmology,'' \textit{EPJC} 85, 100 (2025).
\bibitem{Hui2021} Hui, L., et al., ``Ultralight Scalar Dark Matter,'' \textit{Ann.~Rev.~Astron.~Astrophys.} 59, 247 (2021).
\bibitem{Planck2020} Planck Collaboration, ``Planck 2018 Results VI: Cosmological Parameters,'' \textit{A\&A} 641, A6 (2020).
\bibitem{ADM1959} Arnowitt, R., Deser, S., \& Misner, C.~W. (1959). \textit{Phys.~Rev.} 116, 1322.
\bibitem{MTW1973} Misner, C.~W., Thorne, K.~S., \& Wheeler, J.~A. (1973). \textit{Gravitation.} W.~H.~Freeman.

\end{thebibliography}

\end{document}